\documentclass[11pt,a4paper]{article}
\usepackage[utf8]{inputenc}
\usepackage[vietnamese]{babel}
\usepackage{amsmath,amssymb}
\usepackage{geometry}
\geometry{a4paper, total={180mm,260mm}, left=15mm, right=15mm, top=18mm, bottom=18mm}
\usepackage{tikz}
\usetikzlibrary{shapes.geometric, arrows.meta, positioning, calc}
\usepackage{tcolorbox}
\tcbuselibrary{skins,breakable}
\usepackage{xcolor}
\usepackage{multicol}
\usepackage{array}
\usepackage{fontawesome}
\usepackage{fancyhdr}

\pagestyle{fancy}
\fancyhf{}
\renewcommand{\headrulewidth}{0.5pt}
\renewcommand{\footrulewidth}{0.5pt}
\fancyhead[L]{\small\bfseries\color{teal}\faBook\ SỔ TAY KIẾN THỨC TOÁN 3}
\fancyhead[R]{\small\bfseries\color{teal}Lớp Toán Thầy Lĩnh -- loptoanthaylinh.vn}
\fancyfoot[C]{\small\color{gray}-- Sổ tay kiến thức: Bài 38. Biểu thức số. Tính giá trị của biểu thức số --}

\definecolor{primary}{HTML}{005F73}
\definecolor{secondary}{HTML}{0A9396}
\definecolor{accent}{HTML}{EE9B00}
\definecolor{danger}{HTML}{AE2012}
\definecolor{boxbg}{HTML}{F4F6F8}

\begin{document}

% --- BANNER TIÊU ĐỀ ---
\begin{tcolorbox}[
  enhanced,
  colback=white,
  colframe=primary,
  arc=3mm,
  boxrule=1.2pt,
  left=4mm, right=4mm, top=2.5mm, bottom=2.5mm,
  drop shadow=black!10!white
]
  \centering
  {\Large\bfseries\color{primary}\faBookmark\ BÀI 38: BIỂU THỨC SỐ. TÍNH GIÁ TRỊ CỦA BIỂU THỨC SỐ}\\[1.5mm]
  {\small\itshape\color{black!75}Quy tắc ưu tiên thứ tự thực hiện phép tính: Nhân chia trước, cộng trừ sau, trong ngoặc trước trên 1 trang A4}
\end{tcolorbox}

\vspace{-2mm}


\noindent
\begin{minipage}[t]{0.49\textwidth}\vspace{0pt}
\begin{tcolorbox}[
  enhanced,
  colback=blue!4!white,
  colframe=blue!75!black,
  arc=2mm, boxrule=0.8pt,
  left=2mm, right=2mm, top=1.8mm, bottom=1.8mm,
  title={\footnotesize\bfseries\color{white}\faKey\ 1. Khái niệm Biểu thức số}
]
{\footnotesize
\begin{itemize}\setlength{\itemsep}{1.5pt}\setlength{\leftskip}{-2.5mm}
  \item Các số được nối với nhau bởi các dấu phép tính ($+, -, \times, :$) tạo thành một \textbf{biểu thức số}.
  \item Giá trị của biểu thức số là kết quả nhận được sau khi thực hiện tất cả các phép tính trong biểu thức.
  \item \textbf{Ví dụ:}
  \begin{itemize}\setlength{\itemsep}{0.5pt}\setlength{\leftskip}{-2.5mm}
    \item Biểu thức $125 + 75$ có giá trị là $200$.
    \item Biểu thức $45 \times 2 - 10$ có giá trị là $80$.
  \end{itemize}
\end{itemize}
}
\end{tcolorbox}
\end{minipage}
\hfill
\begin{minipage}[t]{0.49\textwidth}\vspace{0pt}
\begin{tcolorbox}[
  enhanced,
  colback=teal!4!white,
  colframe=teal!75!black,
  arc=2mm, boxrule=0.8pt,
  left=2mm, right=2mm, top=1.8mm, bottom=1.8mm,
  title={\footnotesize\bfseries\color{white}\faSliders\ 2. Ba quy tắc thứ tự phép tính}
]
{\footnotesize
\begin{itemize}\setlength{\itemsep}{1.5pt}\setlength{\leftskip}{-2.5mm}
  \item \textbf{Quy tắc 1 (Chỉ có cộng, trừ HOẶC chỉ có nhân, chia):}
  Thực hiện lần lượt từ \textbf{trái sang phải}.
  \item \textbf{Quy tắc 2 (Có cả nhân, chia VÀ cộng, trừ):}
  Thực hiện các phép tính \textbf{nhân, chia trước}; rồi thực hiện \textbf{cộng, trừ sau}.
  \item \textbf{Quy tắc 3 (Có dấu ngoặc tròn):}
  Thực hiện các phép tính \textbf{trong ngoặc trước}, ngoài ngoặc sau.
\end{itemize}
}
\end{tcolorbox}
\end{minipage}

\vspace{-2.5mm}
% --- KHỐI 3: SƠ ĐỒ TƯ DUY TIKZ ---
\noindent
\begin{tcolorbox}[
  enhanced,
  colback=orange!4!white,
  colframe=orange!80!black,
  arc=2mm, boxrule=0.8pt,
  left=2mm, right=2mm, top=1.5mm, bottom=1.5mm,
  title={\footnotesize\bfseries\color{white}\faSitemap\ SƠ ĐỒ TƯ DUY: CÂY QUYẾT ĐỊNH THỨ TỰ TÍNH GIÁ TRỊ BIỂU THỨC}
]
\centering
\begin{tikzpicture}[
  box/.style={draw, rounded corners=2pt, align=center, font=\scriptsize\bfseries, inner sep=2.5pt, minimum height=0.75cm},
  arr/.style={-{Stealth[length=2mm]}, thick}
]
  \node[box, fill=purple!15, draw=purple!80!black, text width=2.6cm] (b1) at (0, 0) {CÓ NGOẶC $( )$ ?\\Ưu tiên số 1\\Tính trong ngoặc trước};
  \node[box, fill=teal!15, draw=teal!80!black, text width=2.6cm] (b2) at (3.5, 0) {PHÉP NHÂN / CHIA\\Ưu tiên số 2\\Tính trước phép cộng trừ};
  \node[box, fill=blue!12, draw=blue!80!black, text width=2.6cm] (b3) at (7.0, 0) {PHÉP CỘNG / TRỪ\\Ưu tiên số 3\\Tính sau cùng};
  \node[box, fill=orange!15, draw=orange!80!black, text width=2.6cm] (b4) at (10.5, 0) {CÙNG CẤP ĐỘ\\Thực hiện từ\\TRÁI SANG PHẢI};

  \draw[arr, purple!80!black] (b1) -- (b2);
  \draw[arr, teal!80!black] (b2) -- (b3);
  \draw[arr, blue!80!black] (b3) -- (b4);
\end{tikzpicture}
\end{tcolorbox}

\vspace{-2.5mm}
% --- KHỐI 4: VÍ DỤ & MẸO TRÁNH BẪY ---
\noindent
\begin{minipage}[t]{0.49\textwidth}\vspace{0pt}
\begin{tcolorbox}[
  enhanced,
  colback=green!4!white,
  colframe=green!60!black,
  arc=2mm, boxrule=0.8pt,
  left=2mm, right=2mm, top=1.8mm, bottom=1.8mm,
  title={\footnotesize\bfseries\color{white}\faCheckCircle\ Ví dụ thực hành chuẩn mực}
]
{\footnotesize
\textbf{1. Biểu thức không có ngoặc:}\\
$A = 120 + 30 \times 3 = 120 + 90 = \mathbf{210}$.\\
$B = 84 - 24 : 4 = 84 - 6 = \mathbf{78}$.\\
$C = 48 : 2 \times 3 = 24 \times 3 = \mathbf{72}$ (tính từ trái sang phải).\\
\textbf{2. Biểu thức có dấu ngoặc:}\\
$D = (120 + 30) \times 3 = 150 \times 3 = \mathbf{450}$.\\
$E = 72 : (2 \times 4) = 72 : 8 = \mathbf{9}$.
}
\end{tcolorbox}
\end{minipage}
\hfill
\begin{minipage}[t]{0.49\textwidth}\vspace{0pt}
\begin{tcolorbox}[
  enhanced,
  colback=red!4!white,
  colframe=red!70!black,
  arc=2mm, boxrule=0.8pt,
  left=2mm, right=2mm, top=1.8mm, bottom=1.8mm,
  title={\footnotesize\bfseries\color{white}\faExclamationTriangle\ Mẹo tránh bẫy trắc nghiệm}
]
{\footnotesize
\begin{itemize}\setlength{\itemsep}{1.5pt}\setlength{\leftskip}{-2.5mm}
  \item \textbf{Bẫy biểu thức chỉ có nhân và chia:}\\
  Biểu thức $16 : 4 \times 2$, nhiều bạn làm phép nhân $4 \times 2 = 8$ trước rồi lấy $16 : 8 = 2$ là \textbf{SAI}! Cùng cấp nhân chia phải làm từ \textbf{trái sang phải}: $16 : 4 = 4 \implies 4 \times 2 = \mathbf{8}$.
  \item Bẫy tương tự với cộng trừ: $50 - 20 + 10 = 30 + 10 = \mathbf{40}$ (không phải $50 - 30 = 20$).
\end{itemize}
}
\end{tcolorbox}
\end{minipage}


\end{document}
